\section{Introduzione}

La calibrazione intrinseca della camera è il processo con cui si stimano
i parametri interni dell'ottica e del sensore: la lunghezza focale, le
coordinate del punto principale e i coefficienti di distorsione radiale
e tangenziale. Senza questa informazione è impossibile estrarre misure
metriche coerenti dal piano immagine, né ricostruire correttamente la
geometria della scena osservata.

Il progetto si colloca in questo ambito e propone un
flusso di calibrazione che si distingue dall'approccio classico di
Zhang lungo due dimensioni:

\begin{enumerate}[label=\arabic*.]
  \item \textbf{Pattern planare custom.} Invece della scacchiera
        bianco-nera, il tool accetta qualsiasi immagine planare fornita
        dall'utente. I keypoint affini vengono individuati tramite il
        modulo \texttt{cv2.ccalib.CustomPattern} di
        OpenCV, che localizza punti stabili
        sull'immagine di riferimento e li insegue nelle fotografie di
        calibrazione tramite matching di descrittori.

  \item \textbf{Ottimizzazione tramite Algoritmo Genetico (GA).}
        Dopo la prima stima di \(\mat{K}\) con il solver di OpenCV,
        un GA evolve una popolazione di \emph{cromosomi}, ciascuno dei
        quali codifica un sottoinsieme di immagini di calibrazione e un
        insieme di flag numerici. La funzione di fitness è il
        \emph{reprojection error} medio RMSE: il GA minimizza questa
        quantità selezionando il sottoinsieme ottimale di immagini e la
        combinazione ottimale di flag.
\end{enumerate}

\newpage